\documentclass[10pt,a4paper]{article}

\usepackage[a4paper,margin=18mm]{geometry}
\usepackage[T1]{fontenc}
\usepackage[utf8]{inputenc}
\usepackage{array}
\usepackage{booktabs}
\usepackage{enumitem}
\usepackage{longtable}
\usepackage{tabularx}
\usepackage{xcolor}
\usepackage{xurl}
\usepackage[hidelinks]{hyperref}

\definecolor{statusred}{RGB}{150,35,35}
\definecolor{statusamber}{RGB}{130,90,20}
\definecolor{softgray}{RGB}{245,245,245}

\setlength{\parindent}{0pt}
\setlength{\parskip}{0.45em}
\setlength{\tabcolsep}{3pt}
\sloppy
\hbadness=10000
\setlist[itemize]{leftmargin=1.4em,itemsep=0.2em,topsep=0.2em}
\setlist[enumerate]{leftmargin=1.6em,itemsep=0.2em,topsep=0.2em}

\newcommand{\code}[1]{\texttt{\detokenize{#1}}}
\newcommand{\checkbox}{\texttt{[ ]}}
\newcommand{\decision}[1]{\textcolor{statusamber}{\textbf{#1}}}
\newcommand{\blocked}[1]{\textcolor{statusred}{\textbf{#1}}}
\newcolumntype{Y}{>{\raggedright\arraybackslash}X}
\newcolumntype{L}[1]{>{\raggedright\arraybackslash}p{#1}}

\title{\textbf{CD-A OGS Inversion Blocker Discussion Packet}}
\author{Prepared as a self-contained discussion and request packet}
\date{Generated 2026-06-01}

\begin{document}
\maketitle

\section*{Purpose}

This packet puts the remaining external questions, provider file requests, and
modelling-team decisions in one place for discussion. It is intended to be
sent as a standalone document: all questions, required evidence, context, and
response instructions needed by the recipients are included here.

\textbf{Important email-safety note:} this document does not send email. The
mailbox state checked on \code{2026-06-01T03:55:46+02:00} recorded six prepared
request emails, all still unsent. The check found zero sent copies of the
generated request subjects, zero provider-reply hits, and zero recent non-draft
CD-A/HERMES/TeamBeam hits in the searched scope.

\textbf{How recipients should respond:} reply by email with the request id(s) in
the subject, body, or filename. Attach machine-readable source tables, scripts,
calibration notes, geometry files, or provenance notes directly to the reply
where possible. For every response, include source contact, response date,
units, coordinate/reference frame, processing provenance, quality/status flags,
and uncertainty/covariance assumptions where applicable. No access to local
folders or previous file paths is required.

\section*{Current Status Snapshot}

\begin{tabularx}{\textwidth}{@{}L{0.34\textwidth}Y@{}}
\toprule
Item & Current state \\
\midrule
Overall objective readiness & \code{not_complete}. The non-complete requirements are the permeability field/final combined inversion and remaining open questions. \\
Final promotion decision & \blocked{\code{do_not_promote_current_field}}. The current field is an active-objective incumbent only. \\
Open blockers & 8 total: 7 measurement/provider gates and 1 model-provenance confirmation. \\
Request email state & 6/6 expected request emails prepared, 6 still unsent, no sent-subject hits, no provider-reply hits. \\
Recovered material state & Email, file-transfer, and download recovery scans found zero additional gate-closing rows. \\
Report and catalogue state & Active report comments resolved; citation locators pass; measurement traceability passes; frozen-model inclusion audit passes. \\
\bottomrule
\end{tabularx}

\section*{High-Level Discussion Decisions}

\begin{enumerate}
\item \checkbox\ Decide whether to send the six request emails as written, edit them first, or route them differently.
\item \checkbox\ Decide whether to wait for provider responses or explicitly mark some gated streams as diagnostic-only/excluded for the final objective.
\item \checkbox\ Decide the final NMR residual semantics.
\item \checkbox\ Decide the direct-permeability likelihood/support/outlier policy.
\item \checkbox\ After responses or exclusions are recorded, rerun the stream gates, conditional scenarios, field-selection audit, final-promotion checklist, final objective decision register, and objective-readiness audit.
\end{enumerate}

\section{Open Blocker Register}

\begin{longtable}{@{}L{0.19\textwidth}L{0.16\textwidth}L{0.16\textwidth}L{0.21\textwidth}L{0.20\textwidth}@{}}
\toprule
Request id & Priority / stream & Email route & Missing evidence & What closes the gate \\
\midrule
\endhead
\path{ert_transform_support} & high / ERT resistivity & Request email 1; To Gesa; Cc Markus Furche & ERT-to-OGS transform and exact near-niche support mask are assumed, not confirmed. & Coordinate-frame note or transform script, coordinate reference, accepted tunnel/niche support polygon or mask, and decision on 35 cm rock-depth handling. \\
\path{ert_uncertainty} & high / ERT resistivity & Request email 1 & No defensible ERT residual sigma/correlation model; dense VTK pixels cannot be treated as independent rows. & Per-cell, regional, or time-level uncertainty/covariance export, or a simplified weighting rule with log base, filtering, spatial correlation, and effective degrees of freedom. \\
\path{hm_numeric_exports} & high / other HM monitoring & Request email 2; To Gesa & Available collected files contain layout/qualitative evidence, but no hard-residual-ready Geoscope time series, laser statistical product, or full levelling table. & Machine-readable Geoscope mini-piezometer, extensometer, crackmeter, laser-scan interpretation, and levelling tables with time/epoch and support. \\
\path{hm_uncertainty} & high / other HM monitoring & Request email 2 & Metadata for hard HM residual weights are absent. & Unit/reference convention, uncertainty model, covariance/independence assumptions, quality flags, failure periods, laser registration uncertainty/masks, and levelling covariance frame. \\
\path{rh_active_curve_provenance} & high / RH/suction & Request email 3; To Gesa & Active boundary-curve generation, sensor screening, time axis, constants, and extension policy are not confirmed. & Source table/script, sensor selection note, Kelvin conversion constants, model-time origin/timezone, smoothing/manual-edit policy, sign convention, and open/closed mapping. \\
\path{taupe_unit_calibration} & high / Taupe/TDR & Request email 4; To Gesa & Taupe values are not confirmed as volumetric water-content percent, permittivity, or another ARDP-derived proxy. & Workbook unit definition, sensor/band calibration equations, uncertainty by sensor/band, baseline/reference date, and dielectric/water-content conversion notes. \\
\path{perm_endpoint_geometry} & medium / permeability pulse tests & Request email 5; To Gesa & Historical permeability endpoints are needed only if older rows should enter the fit. & Endpoint coordinates, digitized traces, or interval geometry tables for BCD-A24, BCD-A25, BCD-A26, BCD-A27, and BFM-D19 intervals. \\
\path{cte_value_confirmation} & high / model provenance & Request email 6; To Gesa; Cc Tuanny & The XML binds \code{CTE} to solid thermal expansivity, but the value/comment look like a copied heat-capacity entry. & Provider confirmation that the value is intended, inactive/irrelevant, a typo not to be changed in the frozen model, or another documented provenance case. \\
\bottomrule
\end{longtable}

\section{Provider Request Email Pack}

\subsection*{Request Email 1: ERT transform, support, and uncertainty}
\begin{tabularx}{\textwidth}{@{}L{0.24\textwidth}Y@{}}
\toprule
Subject & CD-A ERT transform, support, and uncertainty needed for OGS comparison \\
To / Cc & \code{Gesa.Ziefle@bgr.de} / \code{Markus.Furche@bgr.de} \\
Contact route & Coordinator with named provider Cc. Markus Furche is found in ERT/meeting context, but no direct Gmail sender messages from Markus were found in the scan. \\
\bottomrule
\end{tabularx}

\textbf{Questions to ask / files requested}
\begin{itemize}
\item Confirm the ERT coordinate origin, axes, units, and the currently assumed transform \code{model_x=raw_x}, \code{model_y=raw_y+500}.
\item Provide or confirm the accepted near-niche support mask/polygon used for comparing ERT to OGS.
\item State how 35 cm rock-depth handling and radial support variants should be treated.
\item Provide the ERT uncertainty model: sigma units/log base, time grouping, spatial correlation, filtering rules, unstable-cell handling, and whether cells may be aggregated before comparison.
\end{itemize}

\textbf{Acceptance evidence}
\begin{itemize}
\item ERT spatial projection/support sensitivity can be regenerated without a provisional-transform caveat.
\item ERT can be weighted without treating dense VTK cells as independent residual rows.
\end{itemize}

\textbf{Requested response package:} reply by email with request id
\code{ert_transform_support} and/or \code{ert_uncertainty}. Attach source
tables, scripts/notebooks, polygons/masks, coordinate-frame notes, and
uncertainty/covariance descriptions directly to the reply.

\subsection*{Request Email 2: Geoscope, laser-scan, levelling, and other HM exports}
\begin{tabularx}{\textwidth}{@{}L{0.24\textwidth}Y@{}}
\toprule
Subject & CD-A Geoscope, laser-scan, and levelling numeric exports needed for model validation \\
To / Cc & \code{Gesa.Ziefle@bgr.de} / none \\
Contact route & Coordinator only. Ask Gesa to forward to the relevant BGR Geoscope, laser-scan, and levelling data owners if needed. \\
\bottomrule
\end{tabularx}

\textbf{Questions to ask / files requested}
\begin{itemize}
\item Machine-readable Geoscope mini-piezometer time series.
\item Extensometer and crackmeter time series if available.
\item Laser-scan statistical interpretation products, not only surface geometry.
\item Full precision-levelling tables, including survey epochs and reference convention.
\item For every numeric export: instrument id, timestamp/epoch, value, unit, coordinate/support id, reference/zero convention, processing provenance, quality/status flag.
\item For every candidate residual: uncertainty/covariance or an accepted simplified weighting rule, including failed-sensor and maintenance intervals.
\end{itemize}

\textbf{Acceptance evidence}
\begin{itemize}
\item At least one other-HM request class becomes hard-residual-ready with time/epoch and model-facing support.
\item Each residual has enough metadata to state what OGS quantity/support it measures and what residual weight it receives.
\end{itemize}

\textbf{Requested response package:} reply by email with request id
\code{hm_numeric_exports} and/or \code{hm_uncertainty}. Attach
machine-readable tables and metadata directly to the reply, preferably one file
per instrument class plus a short provenance/readme note.

\subsection*{Request Email 3: RH/suction boundary-curve provenance}
\begin{tabularx}{\textwidth}{@{}L{0.24\textwidth}Y@{}}
\toprule
Subject & CD-A RH/suction boundary-curve provenance needed for OGS forcing \\
To / Cc & \code{Gesa.Ziefle@bgr.de} / none \\
Contact route & Coordinator only; no separate RH/OGS pressure-curve processing owner was identified in the scan. \\
\bottomrule
\end{tabularx}

\textbf{Questions to ask / files requested}
\begin{itemize}
\item Active open-niche pressure-curve source table and any script/notebook used to generate the XML curve.
\item Sensor selection/screening note, including failed or excluded sensors.
\item Model-time zero, timezone, RH percent/fraction convention, pressure unit/sign convention, and temperature/density constants.
\item Kelvin conversion constants and any smoothing/manual-edit or extrapolation policy.
\item Open/closed curve mapping and post-active-curve extension policy.
\end{itemize}

\textbf{Acceptance evidence}
\begin{itemize}
\item The active XML pressure curve can be regenerated or explained well enough that replacement curves/retention targets are no longer based on unknown provenance.
\end{itemize}

\textbf{Requested response package:} reply by email with request id
\code{rh_active_curve_provenance}. Attach the source table, generation
script/notebook, screening note, and any pressure/RH conversion assumptions
directly to the reply.

\subsection*{Request Email 4: Taupe/TDR units and calibration}
\begin{tabularx}{\textwidth}{@{}L{0.24\textwidth}Y@{}}
\toprule
Subject & CD-A Taupe/TDR workbook units and calibration needed before objective activation \\
To / Cc & \code{Gesa.Ziefle@bgr.de} / none \\
Contact route & Coordinator only. Ask Gesa to forward to the Taupe/TDR provider or confirm the source directly. \\
\bottomrule
\end{tabularx}

\textbf{Questions to ask / files requested}
\begin{itemize}
\item Define the units and meaning of \code{Taupe_WC.xlsx} values by sheet/column.
\item Confirm whether values are volumetric water-content percent, permittivity, or another ARDP-derived proxy.
\item Provide sensor/band calibration equations, constants, baseline/reference date, and uncertainty model.
\item State whether matrix/porosity corrections have already been applied.
\end{itemize}

\textbf{Acceptance evidence}
\begin{itemize}
\item The Taupe semantics audit can decide whether absolute Taupe values are water content, permittivity/proxy, or trend-only evidence.
\end{itemize}

\textbf{Requested response package:} reply by email with request id
\code{taupe_unit_calibration}. Attach the workbook unit definition,
calibration equations/constants, baseline/reference information, uncertainty
model, and any matrix/porosity correction notes directly to the reply.

\subsection*{Request Email 5: Historical permeability endpoint geometry}
\begin{tabularx}{\textwidth}{@{}L{0.24\textwidth}Y@{}}
\toprule
Subject & CD-A historical permeability endpoint geometry needed for inactive intervals \\
To / Cc & \code{Gesa.Ziefle@bgr.de} / none \\
Contact route & Coordinator only. Ask Gesa to forward to the BGR permeability data owner if interval endpoints come from another source. \\
\bottomrule
\end{tabularx}

\textbf{Questions to ask / files requested}
\begin{itemize}
\item Endpoint coordinates, labelled digitized traces, or interval geometry tables for historical BCD-A24, BCD-A25, BCD-A26, BCD-A27, and BFM-D19 permeability intervals.
\item For every interval: borehole id, open/closed assignment, start/end coordinates or depths, orientation convention, interval length/support, source date/table, permeability value, uncertainty/evaluation note.
\item If endpoint geometry will not be used, confirm whether these older rows should remain inactive/excluded from the final fit.
\end{itemize}

\textbf{Acceptance evidence}
\begin{itemize}
\item Inactive historical rows can be projected to OGS cells with a trace/support definition rather than excluded for missing geometry.
\end{itemize}

\textbf{Requested response package:} reply by email with request id
\code{perm_endpoint_geometry}. Attach endpoint coordinate tables, digitized
traces, interval-geometry files, or a written confirmation that the older rows
should remain inactive/excluded from the final fit.

\subsection*{Request Email 6: CTE value/provenance confirmation}
\begin{tabularx}{\textwidth}{@{}L{0.24\textwidth}Y@{}}
\toprule
Subject & CD-A OGS CTE value/provenance confirmation needed before thermal-parameter interpretation \\
To / Cc & \code{Gesa.Ziefle@bgr.de} / \code{Tuanny.Cajuhi@bgr.de} \\
Contact route & Coordinator with model-setup Cc. \\
\bottomrule
\end{tabularx}

\textbf{Questions to ask / files requested}
\begin{itemize}
\item Confirm whether the XML \code{CTE} value bound to solid thermal expansivity is intended.
\item If not intended, state whether it is inactive/irrelevant for the current run, a known typo not to be changed in the frozen model, or copied from another parameter such as heat capacity.
\item Provide the provenance/source document or model-setup note for the value.
\end{itemize}

\textbf{Acceptance evidence}
\begin{itemize}
\item The report open-comment audit can move the CTE row from confirmation-required to confirmed inactive, confirmed typo, or confirmed value with the provider response cited.
\end{itemize}

\textbf{Requested response package:} reply by email with request id
\code{cte_value_confirmation}. Include the intended interpretation of the
current XML value and attach or cite the source/provenance note for that value.

\section{Internal Modelling Decisions}

\subsection*{NMR final residual semantics}

\decision{Decision required before final promotion.} The current report keeps
\path{retain_raw_absolute_theta_current_report_default_with_caveats}. The
recommended provisional candidate is \code{within_label_trend_anomaly}, but the
acceptance template records 0/1 approvals.

\begin{tabularx}{\textwidth}{@{}L{0.23\textwidth}L{0.22\textwidth}Y@{}}
\toprule
Option & Current status & Discussion consequence \\
\midrule
\path{raw_absolute_current} & active but conditional & Keeps the current active objective, but compares total NMR-visible water content to mobile OGS water content without a bound/interlayer-water correction. \\
\path{global_uniform_offset} & stress test only & Reject unless collaborators provide one physically justified global offset. \\
\path{absolute_with_label_bias} & conditional & Preserves absolute information but needs approved bias/offset priors and wording that biases are not transport water. \\
\path{within_label_trend_anomaly} & recommended after acceptance & Safest provisional NMR residual because constant bound/interlayer/campaign offsets cancel; absolute water content is not fitted. \\
\bottomrule
\end{tabularx}

\textbf{Required signoff fields:} approval status, primary policy selection,
approver name, approval date, approval reference, and accepted residual
definition.

\subsection*{Direct-permeability likelihood/support policy}

\decision{Decision required before reopening same-support OGS spending or final
field promotion.} The current-report policy is
\path{keep_rowwise_gaussian_default}. The same-support active-objective batch
gate is closed because the current field is already at the support lower bound
for the current support map.

\begin{tabularx}{\textwidth}{@{}L{0.25\textwidth}L{0.23\textwidth}Y@{}}
\toprule
Option & Diagnostic value & Discussion consequence \\
\midrule
\path{keep_rowwise_gaussian_default} & 269.818057 & Reproducible current default, but top row losses and support conflicts must be explicitly accepted. \\
\path{use_robust_row_likelihood} & Student-t diagnostic 121.169 & Needs a written kernel, scale, outlier treatment, and comparison rule against historical row-Gaussian objectives. \\
\path{aggregate_by_model_support_cell} & weighted-mean diagnostic \(1.77\times10^{-28}\) & Needs grouping key, statistic, group weights, and reporting of within-support observation ranges. \\
\path{gate_configured_scalar_outliers} & inside-only 245.675 & Needs row-specific outlier disposition and explicit bounds/tensor-shape decision. \\
\path{new_parameterization_before_more_ogs} & gate option & Requires a named field family and materiality threshold before spending more OGS runs. \\
\bottomrule
\end{tabularx}

\textbf{Support-conflict evidence to acknowledge:} 30 active support cells, 24
repeated support cells, 16 cells spanning at least two log10 units, with the top
conflict at cell 4648, BCD-A32, 0.85-0.87 m, observed range 6.94885 log10.

\subsection*{Final objective scenario selection}

The current field wins only the current raw-NMR scenario and ties the direct-only
scenario. Every currently scored scenario that includes promoted NMR, ERT, or
Taupe selects a different field.

\begin{longtable}{@{}L{0.24\textwidth}L{0.22\textwidth}L{0.20\textwidth}L{0.25\textwidth}@{}}
\toprule
Scenario & Winner & Current rank & What is still required \\
\midrule
\endhead
\path{F01_current_raw_nmr_exclude_gated_streams} & Current field & 1 & Explicitly exclude or keep diagnostic-only all gated streams. \\
\path{F02_promoted_nmr_trend_only} & \path{broad_continuous_001_003_length_0p021m} & 14 & Approve NMR trend/anomaly policy. \\
\path{F03_raw_nmr_plus_ert} & \path{local_basis_sampler_003_basis_029_det_l_0p0100_s_1p000} & 3 & Close ERT transform/support and uncertainty gates. \\
\path{F04_raw_nmr_plus_taupe} & \path{local_bracketed_003_length_0p031m} & 21 & Close Taupe/TDR unit/calibration gate. \\
\path{F05_raw_nmr_plus_ert_taupe} & \path{production_sampler_round3_002_length_0p028m_shift_1p050} & 19 & Close ERT and Taupe/TDR gates. \\
\path{F06_promoted_nmr_plus_ert_taupe} & \path{production_sampler_round3_002_length_0p028m_shift_1p050} & 18 & Approve NMR policy and close ERT/Taupe gates. \\
\path{F07_rank_consensus_diagnostic_selection} & \path{local_basis_sampler_003_basis_029_det_l_0p0100_s_1p000} & 2 & Record explicit waiver that diagnostic ranks, not likelihood, select the field. \\
\path{F08_exclude_nmr_direct_permeability_only} & Current field tie & 1 & Explicitly exclude NMR and decide direct-only tie-break. \\
\path{F09_wait_for_rh_other_hm_endpoint_expansion} & not scored & n/a & Wait for RH, other-HM, or endpoint evidence, then rebuild scenarios. \\
\bottomrule
\end{longtable}

\section{After Replies Arrive: Internal Follow-Up}

\begin{enumerate}
\item Record which request ids were sent, which people received them, and which
answers/files came back.
\item For every received answer, record source contact, response date, file
names, units, coordinate/reference frame, processing provenance, quality/status
flags, uncertainty/covariance assumptions, and which acceptance evidence it
satisfies.
\item Save received files into the project archive under the matching
measurement stream, preserving original filenames where possible and adding a
short readme/provenance note.
\item Rerun the stream gates, conditional scenario matrix, field-selection
audit, final-promotion checklist, final objective decision register, and
objective-readiness audit in the working repository.
\item If new evidence changes the modelling interpretation, update the report
text, citations, measurement traceability notes, and frozen-model inclusion
audit before making any final field-promotion claim.
\end{enumerate}

\section{Packet Provenance}

This document is a self-contained summary of the current CD-A OGS inversion
blockers, provider questions, and modelling-team decisions. It intentionally
does not require recipients to open local folders, internal tracker files, or
previous email drafts. The request ids in this packet are the stable labels to
use when discussing or replying to the open issues.

\end{document}
